\section{Escape velocity results} \label{sec:results}

\begin{figure*}
    \centering
\includegraphics[width=\textwidth]{media/modelcomp3panel_r_8.0_9.0_vmin_310.0_verrcut_0.01_N_5000.png}
    \caption{Comparison of fits to high-speed distribution of stars using the 1PL, 2PL and SEPL models, shown for a representative radial bin of $8-9\,\rm{kpc}$ and at a $\vmin$ of $310\,\rm{km\,s^{-1}}$. The fits at all Galactocentric radii can be found in the appendix. The arrow in the 1PL panel indicates that the escape velocity and its $68\%$ confidence interval are very large and not visible on these axes. The vertical position of the $\vesc$ markers does not encode any information.
    }
    \label{fig:modelcomp3panel}
\end{figure*}

We now apply the fitting procedure outlined in Sec. \ref{sec:methods} to the high-quality DR3 dataset of Sec. \ref{sec:data}, using the 1PL, 2PL and SEPL models of Sec. \ref{sec:modeling}. We perform the fit for each model in $1 \rm{kpc}$ radial bins from $4-11\,\rm{kpc}$. 

In Figure \ref{fig:modelcomp3panel}, we show the data and the best-fit speed distribution for each model, taking a representative radial bin of $8-9\,\rm{kpc}$ and $\vmin$ of $310\,\rm{km\,s^{-1}}$. The fits at all Galactocentric radii can be found in the Appendix, Fig.~\ref{fig:modelcompALL}. The corner plots corresponding to the fits at $8-9\,\rm{kpc}$ are also available in the Appendix (Figs. \ref{fig:corner1PL}, \ref{fig:corner2PL}, \ref{fig:cornerSEPL}).

In Figure \ref{fig:modelcomp3panel}, it can also be seen that the 1PL model results in a best-fit $\vesc$ which is well beyond the fastest star in the sample. Here the shape of the distribution at lower speeds is dictating the behavior of the high-speed tail, with an unreliable extrapolation of the distribution's shape to speeds where there is no data. The best fits for the 2PL model perform better, since the additional model parameters allow for more flexibility in modeling both the lower-speed and fastest stars. However, the 2PL model can also yield a $\vesc$ estimate far from the data region, relying again on strong assumptions about the shape of the distribution at high speeds. This effect is most pronounced for the power law models at large Galactocentric radius, as seen in Fig.~\ref{fig:modelcompALL}. Finally, the SEPL fits exhibit sharper cutoffs near $\vesc$, while simultaneously describing the steep rise of the distribution close to $\vmin$. The SEPL escape velocity estimate thus relies less on extrapolation of the speed distribution, and is very close to the observed cutoff in the distribution. The advantages of using the SEPL are discussed in more detail in Sec. \ref{sec:sepl_advantages}.

\begin{figure*}
    \centering
    \includegraphics[width=\textwidth]{media/profiles_3panel_vmin310.0_verr0.01_N5000_binwidth1.0.pdf}
    \caption{(Top) The escape velocity profile of the high-quality \Gaia~DR3 sample as obtained via the single power law (1PL), 2-component power law (2PL) and stretched exponential power law (SEPL) models, each at $\vmin=310\,\rm{km\,s^{-1}}$. (Bottom left) The stability of each model's escape velocity estimate with increasing $\vmin$, shown here for the representative radial bin of $7-8\,\rm{kpc}$. (Bottom right) Repeating the analysis with the SEPL model using different power law exponents $\gamma$ at high-speed extreme, including the case in which $\gamma$ is treated as a fit parameter. All bands correspond to $68\%$ confidence intervals.}
    \label{fig:3panelprofiles}
\end{figure*}

These features are also apparent when comparing the escape velocity profile of each model, as in the top panel of Fig.~\ref{fig:3panelprofiles}. The 1PL model confidently overestimates the escape velocities, and increases unphysically with Galactocentric radius from $6-9\,\rm{kpc}$. The 2PL model estimates lower escape velocities and relies less on an unreliable extrapolation beyond the data region, at least at small Galactocentric radius. However, at and beyond the Solar position, the 2PL fit produces similar results to the 1PL model and the escape velocity estimates become poorly constrained. The SEPL fitting typically results in the lowest estimates for the escape velocity, as the distribution can exhibit a sharper cutoff at the edge of the data region. Unlike the profiles of the power law models, the SEPL profile is consistent with monotonically decreasing with radius, which is a necessary feature of a physical escape velocity profile. The SEPL escape velocities are most uncertain when it is challenging to distinguish between bound and outlier distributions as in the $5-6\,\rm{kpc}$ bin, and when statistics degrade as at high Galactocentric radius.  

\subsection{Advantages of the SEPL model}
\label{sec:sepl_advantages}

Using the AIC (Eq. \ref{eq:AIC}) to compare the different models reveals that the SEPL and 2PL are not statistically distinguishable ($|\Delta \rm{AIC}|\lesssim 5$) and that the 1PL is strongly disfavored ($|\Delta \rm{AIC}| \simeq 10-100$). Despite the fact that the SEPL and 2PL have similar measures of goodness-of-fit, there are a number of reasons to favor the SEPL as our fiducial result. We discuss below the advantages of using the SEPL model, but note that in our final results for MW halo properties we will also consider $\vesc$ obtained with 2PL model to illustrate the model dependence.

As discussed above, the SEPL model gives results for $\vesc$ that rely less on extrapolation beyond the data region. This is preferable to the overestimation seen in the power law models (on the order of hundreds of $\rm{km\,s^{-1}}$ at high Galactocentric radius) as it is an approach informed more by the data than assumptions about the tail shape. In light of this discussion, one might ask whether it is appropriate to use the speed of the fastest star as an estimator for $\vesc$. The difficulty of this approach is that it does not account for outliers or statistical fluctuations near the tail, and degrades significantly in low statistics regimes such as at high Galactocentric radius. Using the SEPL model with a cutoff in the distribution accounts for these effects, and results in larger uncertainties on $\vesc$ in radial bins with limited statistics or outliers.

In addition, using the SEPL model mitigates one of the issues seen in power law models: the power law slope that governs the low-speed behavior, $k$, is often both unconverged and correlated with $\vesc$. In the Appendix, Fig.~\ref{fig:correlation} shows the degree to which the model parameters of the 1PL, 2PL and SEPL are correlated with $\vesc$ as a function of Galactocentric radius. In both the 1PL and 2PL models, $k$ exhibits a high correlation coefficient with $\vesc$, and it can be seen in the representative corner plots Fig. \ref{fig:corner1PL} and \ref{fig:corner2PL} that this parameter is unconverged even with a large prior range. This suggests that the marginal distribution for $\vesc$ is in fact not reliable in this case, as it is correlated with an unconverged parameter. We find the $\beta$ parameter in the SEPL model is sometimes similarly unconverged, in part because the upper prior limit is set to avoid numerical overflow. However, it is \textit{uncorrelated with the escape velocity} and thus does not pose the same problem as the power law models. There is one radial bin ($5-6\,\rm{kpc}$) in which there is meaningful correlation between $\beta$ and $\vesc$, but this is accompanied by a large uncertainty on the $\vesc$ estimate of the SEPL, and so does not significantly bias results of DM halo parameter fitting in Sec.~\ref{sec:DM}.

The lower left panel of Fig.~\ref{fig:3panelprofiles} shows the result of increasing $\vmin$ for the three models considered in this paper. For the power law models, the results for $\vesc$ become highly uncertain at high $\vmin$ as the statistics of the sample degrade, with both models estimating an escape velocity $\sim 150\,\rm{km\,s^{-1}}$ above the SEPL and low-$\vmin$ 2PL measurements. 
We find that the SEPL model gives results that are more stable to the choice of $\vmin$ than the power law models. This can be explained by the greater flexibility in the SEPL for describing the rise toward $\vmin$, which better describes the features observed in the data. Using the AIC to compare the models, the SEPL and 2PL results are indistinguishable at all $\vmin$ while all three models (1PL, 2PL, SEPL) become indistinguishable for high $\vmin$, consistent with expectations given the limited statistics. 

We also consider the more general form of the SEPL, Eq.~\ref{eq:generalSEPL}, with various choices of the power law index $\gamma$. The lower right panel of Fig.~\ref{fig:3panelprofiles} shows fit results with different fixed $\gamma=2,3$ as well as with $\gamma$ as a free parameter.  Larger $\gamma$ values give distributions with shallower cutoffs in the tail, which tend to result in higher $\vesc$ values. The SEPL with $\gamma$ a free parameter reproduces the behavior of the 2PL model close to $\vesc$, in particular rising and becoming more uncertain past the solar position. This feature could be due to the reduced number of stars above $\vmin$ in these bins, and the ``overshooting" issue of the power law models discussed in Section~\ref{sec:limitations_2pl}. We use $\gamma=1$ since it results in the lowest values of $\vesc$ in this comparison, consistent with our desire to avoid extrapolating beyond the data region. Further reducing $\gamma$ below 1 to give an even sharper cutoff would not substantially impact our results, since there are limited statistics to distinguish such low values of $\gamma$. For example, comparing the $\gamma=1$ and $\gamma=2$ cases we see these profiles are already largely consistent within $1\sigma$, but with the $\gamma=2$ case systematically above $\gamma=1$.

Lastly, we test the consistency of the escape velocity profile with an alternative radial binning (but the same bin width), offset by $0.5\,\rm{kpc}$ from the fiducial bins. This is shown in the Appendix Fig. \ref{fig:offsetbins} for both 2PL and SEPL models. The offset escape velocity profiles exhibit the same features as the fiducial profiles, confirming that bin membership uncertainty has a small effect on the resulting escape velocity results. 

\begin{figure*}[t]
    \centering
    \includegraphics[width=\textwidth]{media/profiles_litcomp_and_DM_4.0_11.0_vmin310.0_verrcut0.01_cutoff_powerlaw_varscale_fixpowerlaw_NFW.pdf}
    \caption{
    (Top) stretched exponential power law (SEPL) and 2-component power law (2PL) escape velocity profiles along with comparable results from the literature. Bands and error bars correspond to $68\%$ confidence intervals. Horizontal error bars for isolated results represent the Galactocentric radius range of sources used for that result. The escape velocities shown from the literature do not include any corrections motivated by simulations. (Bottom) Best fit of the NFW dark matter halo to the escape velocity profile, using SEPL escape velocity profile and assuming a known baryonic content of the Milky Way. The NFW halo parameters $\rm{M}_{200c}^{\rm{DM}}$ and $c_{200c}^{\rm{DM}}$ represent the virial mass and concentration of the dark matter halo. The SEPL profile is shown at $68\%$ (boxes) and $95\%$ (whiskers) confidence intervals. The baryon-only and NFW-only profiles show the escape velocity if the MW consisted of only those components, respectively. }
    \label{fig:litcomp_and_DM_profiles}
\end{figure*}

\subsection{Comparison with previous work}

For our final results, we consider the 2PL and SEPL models with the fiducial radial bins, and a default power law index of $\gamma =1$ for SEPL. We take $\vmin=310\,\rm{km\,s^{-1}}$ as the default because this maintains a large sample size while still remaining above the typically chosen $300\,\rm{km\,s^{-1}}$ to avoid disk contamination. Furthermore, as discussed above, the power law models perform less well at high $\vmin$ given the lower statistics. 

In Figure \ref{fig:litcomp_and_DM_profiles}, we present the comparison of the 2PL and SEPL escape velocity profiles with individual measurements and profiles from the literature. We find an escape velocity of $486^{+10}_{-5}\,\rm{km\,s^{-1}}$ for the SEPL model and $539^{+69}_{-34}\,\rm{km\,s^{-1}}$ for the 2PL model in the $8-9\,\rm{kpc}$ bin, which is largely consistent with previous work obtaining $\vesc$ near the solar radius. For the \cite{Necib2022} result, the escape velocity is consistent within $1\sigma$ for the SEPL measurement and $\sim 1.5\sigma$ for the 2PL measurement, but the escape velocity found in this work is larger by $\sim20-30\,\rm{km\,s^{-1}}$. This is likely higher due to the larger statistics of \Gaia~DR3 as compared to eDR3, allowing for stricter quality cuts and thus a more reliable discrimination between outlier and bound distributions, in addition to speed distributions which are filled closer to the escape velocity.

Compared to the measurements of \cite{deason:2019} and \cite{Koppelman2021} near the solar radius, our $\vesc$ estimates are consistent. Note that the values cannot be compared directly, since the analysis of \cite{deason:2019} utilizes an ansatz about the shape of the escape velocity profile, and the results of \cite{Koppelman2021} are obtained using a 4 kpc radial bin (horizontal bar), compared to our results in 1 kpc bins. The $\vesc$ parameter is known to be correlated with the fastest star in the distribution being fit, and if using a wide radial bin, one will likely preferentially fit the high-speed stars which are expected to come from the smallest radii. Examining the horizontal error bars of Fig. \ref{fig:litcomp_and_DM_profiles}, it is plausible that the \cite{Koppelman2021} measurement is more appropriately compared with the 2PL or SEPL measurements at the lower radii of $\sim6-7$ kpc.

Turning to the escape velocity profile, the key feature of the SEPL escape velocity profile is that it trends toward lower $\vesc$ at large Galactocentric radii, consistent with expectation. This is in contrast to fits with power law models, which exhibit a significant rise at and beyond the solar position. Our results with the 2PL model are similar to the profile of \cite{Koppelman2021}, which exhibits approximately the same rise in $\vesc$ at large Galactocentric radii. Note that in \cite{Koppelman2021}, in order to obtain what we label as the ``inferred profile" (Fig. \ref{fig:litcomp_and_DM_profiles}), first a fit is performed using stars within $2\,\rm{kpc}$ of the solar position. Then using the posterior for $k$ in this local fit as a {\emph{prior}} on $k$, escape velocity fits are obtained in $0.5\,\rm{kpc}$ bins from $4-11\,\rm{kpc}$, giving an escape velocity profile assuming the power law slope is consistent across Galactocentric radii. We also show the profile of \cite{Monari2018}, which was obtained using \Gaia~DR2 with independent fits to $\vesc$ and $k$ in radial bins, using a bootstrap method to account for uncertainties in bin membership and speed. It does not exhibit a significant trend due to the large uncertainties resulting from the relatively small statistics of \Gaia~DR2, but also gives high $\vesc$ estimates at large Galactocentric radii. These features suggest that power law modeling does not reliably describe the high-speed tail, at least at larger radii.

Note that \cite{williams:2017} also fit an escape velocity profile using data from the Sloan Digital Sky Survey \citep{SDSS:2012}, but the modelling was performed in a distinct manner from this work and the other papers discussed in this section. In particular, \cite{williams:2017} makes an ansatz for the radial dependence of the escape velocity profile, fits the ansatz to data over a wide range of radial values simultaneously, and then uses the result to obtain the escape velocity at the solar position. The paper infers a local escape velocity of $\vesc(r_\odot) = 521^{+46}_{-30}\,{\rm{km\,s^{-1}}}$, consistent with previous estimates at that position \citep{Monari2018, deason:2019, Koppelman2021} and the  SEPL and 2PL inference of this work at the $\sim 1\sigma$ level, and consistent with the local estimates of \cite{Necib2022} at the $\sim 2\sigma$ level. Additionally, in \cite{Prudil:2022} the 1PL modelling approach was applied to speed data from RR Lyrae stars from $4-12\,{\rm{kpc}}$ to obtain a measurement of $\vesc(r_\odot) = 512^{+94}_{-37}\,{\rm{km\,s^{-1}}}$, which is consistent with all measurements discussed here, likely in part due to the large range in Galactocentric radii used for the inference.